\section{Finite windows}
\label{sec:finite_windows}

By \Cref{thm:rh:concrete-window-dc}, a proof along the duality-confinement
route needs upper bounds tending to zero for the window sums \(A_W\).
This section records some elementary facts about natural candidates
for such bounds.  The main point is \Cref{prop:rh:window-identities}(b):
the symmetric probe \(L_W(x)=\sum x_\rho\) sees nothing of the odd
displacement vector, by symmetry alone, so its residual charge is the
whole window sum; the equivalences of
\Cref{prop:rh:window-equivalences} are then tautologies forced by this
choice of probe, and carry no information about \(\zeta\).  Only the oddness of the displacement under the
functional-equation map is used here; no fixed locus is involved.

\paragraph{Setting.}
Throughout, \(\sigma(\rho)=1-\rho\) is the functional-equation map on
the nontrivial zeros (not to be confused with the reflection \(\JL\) in
the critical line), a finite set \(W\) of nontrivial
zeros is \emph{reflected} if \(\sigma(W)=W\), and the positive integer
weights satisfy \(m_{\sigma(\rho)}=m_\rho\).  A sequence \((W_n)\) of
finite sets is \emph{exhaustive} if every nontrivial zero lies in
\(W_n\) for all large \(n\); exhaustive sequences of reflected windows
exist because the nontrivial zeros are countable and each orbit of
\(\sigma\) has at most two elements.  These windows differ from those of
\Cref{thm:rh:concrete-window-dc}, which are closed under
\(\JL(\rho)=1-\bar\rho\) and use unit weights; windows closed under both
reflections also exist, since the two reflections commute and each
orbit of the group they generate has at most four elements.  On the Euclidean
space \(\mathbb R^W\) put
\[
  v_W(\rho)=\sqrt{m_\rho}\Bigl(\Real\rho-\frac12\Bigr),
  \qquad
  L_W(x)=\sum_{\rho\in W}x_\rho ,
\]
the \emph{displacement vector} and the \emph{invariant probe}.  For a linear map \(L\) from \(\mathbb R^W\) to a finite-dimensional
real inner product space, let \(L^+\) be its Moore--Penrose pseudoinverse and
\[
  \Xi(L)=I-L^+L ,
\]
the orthogonal projection onto \(\ker L\); \(L^+L\) is the orthogonal
projection onto \((\ker L)^\perp\).  For a vector \(v\), the
\emph{residual charge} \(\langle v,\Xi(L)v\rangle=\|\Xi(L)v\|^2\) measures
the part of \(v\) that the probe \(L\) does not see.  We call \(LL^*\) the \emph{native currency} of \(L\).  (These are the
identity-audit cases of the residual, native currency and decoder
\(L^+\) of the companion paper \citep{TsiokosXiCompanion2026}.)  A \emph{payment} for a window sum is an upper bound for it,
or for a quadratic form dominating it, that one hopes to derive
independently.  Operator inequalities \(S\preceq T\) are in the sense of
quadratic forms.

\begin{proposition}[Window identities]
\label{prop:rh:window-identities}
Let \(W\) be a reflected window with symmetric positive weights.
\begin{enumerate}[label=\textup{(\alph*)}]
  \item \(\|v_W\|^2=A_W\), and \(A_W\leq\AZ\).
  \item \(L_Wv_W=0\).  Consequently
    \(\langle v_W,\Xi(L_W)v_W\rangle=A_W\): the invariant probe sees
    nothing of the displacement, and the residual charge is the whole
    window sum.
  \item For every linear map \(L\) on \(\mathbb R^W\),
    \(\langle v_W,(L^+(LL^*)(L^+)^*+\Xi(L))v_W\rangle=A_W\).  Changing the
    probe moves charge between the decoded native currency and the
    residual, but never changes their total.
  \item Let \(P_W\) be the orthogonal projection onto \(\ker L_W\), a probe
    determined by the window alone (it does not depend on the
    displacements).  Then
    \(\langle v_W,\Xi(P_W)v_W\rangle=0\) and
    \(\langle v_W,P_WP_W^*v_W\rangle=A_W\).  In particular, a window
    containing a zero off the critical line has zero residual charge for this probe and strictly positive native currency.
\end{enumerate}
Part (d) shows that residual charge depends on the probe; the
equivalences of \Cref{prop:rh:window-equivalences} concern the
invariant probe \(L_W\).
\end{proposition}

\begin{proof}
(a) is the definition of \(v_W\) and of the extended sum.  (b): the
reflection pairs \(\rho\) with \(1-\rho\), which has the opposite
displacement and the same weight; the only fixed point \(\rho=\frac12\)
has displacement zero.  So the terms of \(L_Wv_W\) cancel, \(v_W\) lies in
\(\ker L_W\), and \(\Xi(L_W)v_W=v_W\).  (c): \(L^+(LL^*)(L^+)^*=L^+L\) is
the orthogonal projection onto \((\ker L)^\perp\), so the operator in
(c) is the identity.  (d): \(\ker P_W=(\ker L_W)^\perp\), which is
orthogonal to \(v_W\) by (b); and \(P_WP_W^*=P_W\) fixes \(v_W\).
\end{proof}

\begin{proposition}[Two elementary limitations]
\label{prop:rh:window-obstructions}
Let \(W\) be a nonempty window.
\begin{enumerate}[label=\textup{(\alph*)}]
  \item \emph{Operator transport bounds for the identity target cannot
    vanish.}  This part does not use the zeros at all, and the same
    argument works for any finite-dimensional target of \(L\).  Let
    \(L:\mathbb R^W\to\mathbb R\) be linear, and suppose
    \(LL^*\preceq\Theta\), \(\Xi(L)\preceq\Omega\) and
    \(L^+\Theta(L^+)^*+\Omega\preceq H\cdot I\) for operators
    \(\Theta,\Omega\) and a real number \(H\).  Then \(H\geq1\).
  \item \emph{The probe cannot be coercive on its own residual.}  If
    \(c\|x\|\leq|L_W x|\) for some \(c>0\) and all \(x\) in the range of
    \(\Xi(L_W)\), then \(L_W\) is injective.  Hence no such \(c\) exists when \(W\) has at
    least two elements.
\end{enumerate}
\end{proposition}

\begin{proof}
(a) Conjugating \(LL^*\preceq\Theta\) by \(L^+\) and adding
\(\Xi(L)\preceq\Omega\) gives
\(I=L^+L+\Xi(L)\preceq L^+\Theta(L^+)^*+\Omega\preceq H\cdot I\), and
\(\mathbb R^W\neq0\).  (b) The range of \(\Xi(L_W)\) is \(\ker L_W\), on
which \(L_W\) vanishes; the inequality forces \(\ker L_W=0\).  A nonzero
linear functional on \(\mathbb R^W\) is injective only if
\(\dim\mathbb R^W=1\).
\end{proof}

Here the \emph{identity target} is the map \(I\) on \(\mathbb R^W\) whose
values the payment is meant to control, and a \emph{transport bound} is
the scalar \(H\) in an inequality of the form in (a), which dominates the
decoded native budget plus the residual bound by \(H\cdot I\).  Part (b) says that coercivity on the residual forces \(L_W\) to be
injective, which fails as soon as \(W\) has at least two elements,
because then \(\ker L_W\neq0\).  Part (a) says only that
the identity operator is not small: a transport
bound of the form in (a), which dominates the identity target as an
operator rather than on the actual displacement vector \(v_W\), cannot
tend to zero.

\begin{proposition}[Payments that restate the Riemann hypothesis]
\label{prop:rh:window-equivalences}
Let \((W_n)\) be an exhaustive sequence of reflected windows with
symmetric positive weights, and write \(v_n=v_{W_n}\), \(L_n=L_{W_n}\),
\(A_n=A_{W_n}\).  Each of the following is equivalent to the Riemann
hypothesis:
\begin{enumerate}[label=\textup{(\alph*)}]
  \item \(\langle v_n,\Xi(L_n)v_n\rangle=0\) for every \(n\);
  \item there are \(b_n\to0\) with \(\langle v_n,\Xi(L_n)v_n\rangle\leq b_n\);
  \item there are operators \(\Omega_n\succeq\Xi(L_n)\) and \(b_n\to0\)
    with \(\langle v_n,\Omega_nv_n\rangle\leq b_n\);
  \item for a fixed real number \(q<1\),
    \(\langle v_n,\Xi(L_n)v_n\rangle\leq q\,A_n\) for every \(n\);
  \item \(\sum_n A_n<\infty\).
\end{enumerate}
\end{proposition}

\begin{proof}
By \Cref{prop:rh:window-identities}(b), the residual charge in (a),
(b) and (d) equals \(A_n\), and in (c) it is at most
\(\langle v_n,\Omega_nv_n\rangle\).  If some zero \(\rho_0\) is off the
critical line, it lies in \(W_n\) for all large \(n\), and then
\(A_n\geq(\Real\rho_0-\frac12)^2>0\) for all large \(n\); this
contradicts each of (a)--(e) (for (d), \(A_n\leq qA_n\) forces
\(A_n=0\); for (e), the terms of a convergent series tend to zero).
Conversely, under the Riemann hypothesis every \(A_n\) is zero, and
(a)--(e) hold with \(\Omega_n=\Xi(L_n)\) and \(b_n=0\).
\end{proof}

The proposition does not say that the Riemann hypothesis is hard; it
says that, on actual zero windows, bounding the residual charge,
bounding any operator above it, asking for a uniform residual gap, or
asking for summable window sums is each the Riemann hypothesis in
another form.  An argument along these lines becomes a proof only if
the bound is derived from zeta-function data by a route that does not
already presuppose it.
